\documentclass[a4paper,11pt]{book}
%\documentclass[a4paper,twoside,11pt,titlepage]{book}
\usepackage{listings}
\usepackage[utf8]{inputenc}
\usepackage[spanish]{babel}

% \usepackage[style=list, number=none]{glossary} %
%\usepackage{titlesec}
%\usepackage{pailatino}

\decimalpoint
\usepackage{dcolumn}
\newcolumntype{.}{D{.}{\esperiod}{-1}}
\makeatletter
\addto\shorthandsspanish{\let\esperiod\es@period@code}
\makeatother


%\usepackage[chapter]{algorithm}
\RequirePackage{verbatim}
%\RequirePackage[Glenn]{fncychap}
\usepackage{fancyhdr}
\usepackage{graphicx}
\usepackage{afterpage}

\usepackage{longtable}

\usepackage[pdfborder={000}]{hyperref} %referencia

% ********************************************************************
% Re-usable information
% ********************************************************************
\newcommand{\myTitle}{Deep Learning para diagnóstico a partir de imágenes Biomédicas\xspace}
\newcommand{\myDegree}{Grado en Ingeniería Informática\xspace}
\newcommand{\myName}{Francisco Carrillo Pérez \xspace}
\newcommand{\myProf}{Luis Javier Herrera Maldonado\xspace}
%\newcommand{\myOtherProf}{Nombre Apllido1 Apellido2 (tutor2)\xspace}
%\newcommand{\mySupervisor}{Put name here\xspace}
\newcommand{\myFaculty}{Escuela Técnica Superior de Ingenierías Informática y de
Telecomunicación\xspace}
\newcommand{\myFacultyShort}{E.T.S. de Ingenierías Informática y de
Telecomunicación\xspace}
\newcommand{\myDepartment}{Departamento de Arquitectura de Computadores\xspace}
\newcommand{\myUni}{\protect{Universidad de Granada}\xspace}
\newcommand{\myLocation}{Granada\xspace}
\newcommand{\myTime}{\today\xspace}
\newcommand{\myVersion}{Version 0.1\xspace}


%\hypersetup{
%pdfauthor = {\myName franciscocp@correo.ugr.es},
%pdftitle = {\myTitle},
%pdfsubject = {},
%pdfkeywords = {palabra_clave1, palabra_clave2, palabra_clave3, ...},
%pdfcreator = {LaTeX con el paquete ....},
%pdfproducer = {pdflatex}
%}

%\hyphenation{}


%\usepackage{doxygen/doxygen}
\usepackage{pdfpages}
\usepackage{url}
\usepackage{colortbl,longtable}
\usepackage[stable]{footmisc}
%\usepackage{index}
%\makeindex
%\usepackage[style=long, cols=2,border=plain,toc=true,number=none]{glossary}
%\makeglossary

% Definición de comandos que me son tiles:
%\renewcommand{\indexname}{Índice alfabético}
%\renewcommand{\glossaryname}{Glosario}

\pagestyle{fancy}
\fancyhf{}
\fancyhead[LO]{\leftmark}
\fancyhead[RE]{\rightmark}
\fancyhead[RO,LE]{\textbf{\thepage}}
\renewcommand{\chaptermark}[1]{\markboth{\textbf{#1}}{}}
\renewcommand{\sectionmark}[1]{\markright{\textbf{\thesection. #1}}}

\setlength{\headheight}{1.5\headheight}

\newcommand{\HRule}{\rule{\linewidth}{0.5mm}}
%Definimos los tipos teorema, ejemplo y definición podremos usar estos tipos
%simplemente poniendo \begin{teorema} \end{teorema} ...
\newtheorem{teorema}{Teorema}[chapter]
\newtheorem{ejemplo}{Ejemplo}[chapter]
\newtheorem{definicion}{Definición}[chapter]

\definecolor{gray97}{gray}{.97}
\definecolor{gray75}{gray}{.75}
\definecolor{gray45}{gray}{.45}
\definecolor{gray30}{gray}{.94}

\lstset{ frame=Ltb,
     framerule=0.5pt,
     aboveskip=0.5cm,
     framextopmargin=3pt,
     framexbottommargin=3pt,
     framexleftmargin=0.1cm,
     framesep=0pt,
     rulesep=.4pt,
     backgroundcolor=\color{gray97},
     rulesepcolor=\color{black},
     %
     stringstyle=\ttfamily,
     showstringspaces = false,
     basicstyle=\scriptsize\ttfamily,
     commentstyle=\color{gray45},
     keywordstyle=\bfseries,
     %
     numbers=left,
     numbersep=6pt,
     numberstyle=\tiny,
     numberfirstline = false,
     breaklines=true,
   }
 
% minimizar fragmentado de listados
\lstnewenvironment{listing}[1][]
   {\lstset{#1}\pagebreak[0]}{\pagebreak[0]}

\lstdefinestyle{CodigoC}
   {
	basicstyle=\scriptsize,
	frame=single,
	language=C,
	numbers=left
   }
\lstdefinestyle{CodigoC++}
   {
	basicstyle=\small,
	frame=single,
	backgroundcolor=\color{gray30},
	language=C++,
	numbers=left
   }

 
\lstdefinestyle{Consola}
   {basicstyle=\scriptsize\bf\ttfamily,
    backgroundcolor=\color{gray30},
    frame=single,
    numbers=none
   }


\newcommand{\bigrule}{\titlerule[0.5mm]}


%Para conseguir que en las páginas en blanco no ponga cabecerass
\makeatletter
\def\clearpage{%
  \ifvmode
    \ifnum \@dbltopnum =\m@ne
      \ifdim \pagetotal <\topskip
        \hbox{}
      \fi
    \fi
  \fi
  \newpage
  \thispagestyle{empty}
  \write\m@ne{}
  \vbox{}
  \penalty -\@Mi
}
\makeatother

\usepackage{pdfpages}
\begin{document}
\begin{titlepage}
 
 
\newlength{\centeroffset}
\setlength{\centeroffset}{-0.5\oddsidemargin}
\addtolength{\centeroffset}{0.5\evensidemargin}
\thispagestyle{empty}

\noindent\hspace*{\centeroffset}\begin{minipage}{\textwidth}

\centering
\includegraphics[width=0.9\textwidth]{imagenes/logo_ugr.jpg}\\[1.4cm]

\textsc{ \Large TRABAJO FIN DE GRADO\\[0.2cm]}
\textsc{ INGENIERÍA EN INGENIERÍA INFORMÁTICA}\\[1cm]
% Upper part of the page
% 
% Title
{\Huge\bfseries Deep Learning para diagnóstico a partir de imágenes Biomédicas\\
}
\noindent\rule[-1ex]{\textwidth}{3pt}\\[3.5ex]
{\large\bfseries Subtitulo del Proyecto}
\end{minipage}

\vspace{2.5cm}
\noindent\hspace*{\centeroffset}\begin{minipage}{\textwidth}
\centering

\textbf{Autor}\\ {Francisco Carrillo Pérez}\\[2.5ex]
\textbf{Directores}\\
Luis Javier Herrera Maldonado\\ [2cm]
\includegraphics[width=0.3\textwidth]{imagenes/etsiit_logo.png}\\[0.1cm]
\textsc{Escuela Técnica Superior de Ingenierías Informática y de Telecomunicación}\\
\textsc{---}\\
Granada, mes de 201
\end{minipage}
%\addtolength{\textwidth}{\centeroffset}
%\vspace{\stretch{2}}
\end{titlepage}



\chapter*{}
%\thispagestyle{empty}
%\cleardoublepage

%\thispagestyle{empty}

\begin{titlepage}
 
 
\setlength{\centeroffset}{-0.5\oddsidemargin}
\addtolength{\centeroffset}{0.5\evensidemargin}
\thispagestyle{empty}

\noindent\hspace*{\centeroffset}\begin{minipage}{\textwidth}

\centering
%\includegraphics[width=0.9\textwidth]{imagenes/logo_ugr.jpg}\\[1.4cm]

%\textsc{ \Large PROYECTO FIN DE CARRERA\\[0.2cm]}
%\textsc{ INGENIERÍA EN INFORMÁTICA}\\[1cm]
% Upper part of the page
% 

 \vspace{3.3cm}

%si el proyecto tiene logo poner aquí
\includegraphics{imagenes/logo.png} 
 \vspace{0.5cm}

% Title

{\Huge\bfseries Título del proyecto\\
}
\noindent\rule[-1ex]{\textwidth}{3pt}\\[3.5ex]
{\large\bfseries Subtítulo del proyecto.\\[4cm]}
\end{minipage}

\vspace{2.5cm}
\noindent\hspace*{\centeroffset}\begin{minipage}{\textwidth}
\centering

\textbf{Autor}\\ {Nombre Apellido1 Apellido2 (alumno)}\\[2.5ex]
\textbf{Directores}\\
{Nombre Apellido1 Apellido2 (tutor1)\\
Nombre Apellido1 Apellido2 (tutor2)}\\[2cm]
%\includegraphics[width=0.15\textwidth]{imagenes/tstc.png}\\[0.1cm]
%\textsc{Departamento de Teoría de la Señal, Telemática y Comunicaciones}\\
%\textsc{---}\\
%Granada, mes de 201
\end{minipage}
%\addtolength{\textwidth}{\centeroffset}
\vspace{\stretch{2}}

 
\end{titlepage}






\cleardoublepage
\thispagestyle{empty}

\begin{center}
{\large\bfseries Deep Learning para diagnóstico a partir de imágenes Biomédicas}\\
\end{center}
\begin{center}
Francisco Carrillo Pérez (alumno)\\
\end{center}

%\vspace{0.7cm}
\noindent{\textbf{Palabras clave}:DeepLearning, CNN, Alzheimer}\\

\vspace{0.7cm}
\noindent{\textbf{Resumen}}\\

La enfermedad de Alzheimer es una de las enfermedades que más afecta a pacientes de edad avanzada en todo el mundo. Su cura se desconoce, por lo que un diagnóstico precoz puede ayudar a mejorar notablemente la vida del paciente. El problema es que la clasificación de la enfermedad en edad temprana es una tarea complicada, además de que se puede confundir con otros deterioros que acaecen propios de la edad. Con los nuevos avances que se han obtenido en el uso de técnicas de Deep Learning en el área de Visión por Computador y clasificación de imágenes, el uso de estas técnicas para la clasificación de pacientes puede suponer una ayuda notable a la hora de que se pueda diagnosticar correctamente si un paciente está comenzando a desarrollar esta enfermedad, con lo que se podría tratar con suficiente tiempo. Con este trabajo, se intenta responder si con estas técnicas podemos realizar una clasificación correcta de imágenes 2D cerebrales de pacientes y de ser así cuáles serían las capas del cerebro más favorables a la hora de realizar esta clasificación.
\cleardoublepage


\thispagestyle{empty}


\begin{center}
{\large\bfseries Deep Learning for diagnosis using Biomedical images}\\
\end{center}
\begin{center}
Francisco Carrillo Pérez (student)\\
\end{center}

%\vspace{0.7cm}
\noindent{\textbf{Keywords}: DeepLearning, CNN, Alzheimer}\\

\vspace{0.7cm}
\noindent{\textbf{Abstract}}\\

Alzheimer disease is one of the diseases that affect the most to the elderly in the whole world. The cure is unknown, so an early diagnosis is very important for making improves in the pacient life. The problem is that the diagnosis of the disease in an early stage is really difficult, adding the fact that it could be difficult to distinguish in respect to other damages of age. With the new development of Deep Learning techniques in the area of Computer Vision and Image Classification, the use of this techniques for pacient classification could be a usefull help for when deciding if a pacient is developing the disease, helping to treat it in an early stage. With this TFG, we are trying to anwer if witj this techniques we can make a good classification of 2D brain images of pacientes and if so, which are the best slices of the brain for doing this classification.
\chapter*{}
\thispagestyle{empty}

\noindent\rule[-1ex]{\textwidth}{2pt}\\[4.5ex]

Yo, \textbf{Francisco Carrillo Pérez}, alumno de la titulación Ingeniería Informática de la \textbf{Escuela Técnica Superior
de Ingenierías Informática y de Telecomunicación de la Universidad de Granada}, con 77140580Y, autorizo la
ubicación de la siguiente copia de mi Trabajo Fin de Grado en la biblioteca del centro para que pueda ser
consultada por las personas que lo deseen.

\vspace{6cm}

\noindent Fdo: Francisco Carrillo Pérez

\vspace{2cm}

\begin{flushright}
Granada a X de mes de 201 .
\end{flushright}


\chapter*{}
\thispagestyle{empty}

\noindent\rule[-1ex]{\textwidth}{2pt}\\[4.5ex]

D. \textbf{Luis Javier Herrera Maldonado (tutor1)}, Profesor del Área de XXXX del Departamento YYYY de la Universidad de Granada.

\vspace{0.5cm}

D. \textbf{ALberto Guillén perales (tutor2)}, Profesor del Área de XXXX del Departamento YYYY de la Universidad de Granada.


\vspace{0.5cm}

\textbf{Informan:}

\vspace{0.5cm}

Que el presente trabajo, titulado \textit{\textbf{Deep Learning para diagnóstico a partir de imágenes Biomédicas}},
ha sido realizado bajo su supervisión por \textbf{Francisco Carrillo Pérez (alumno)}, y autorizamos la defensa de dicho trabajo ante el tribunal
que corresponda.

\vspace{0.5cm}

Y para que conste, expiden y firman el presente informe en Granada a X de mes de 201 .

\vspace{1cm}

\textbf{Los directores:}

\vspace{5cm}

\noindent \textbf{Nombre Apellido1 Apellido2 (tutor1) \ \ \ \ \ Nombre Apellido1 Apellido2 (tutor2)}

\chapter*{Agradecimientos}
\thispagestyle{empty}

       \vspace{1cm}


Poner aquí agradecimientos...


\frontmatter
\tableofcontents
\listoffigures
\listoftables
%
%\mainmatter
%\setlength{\parskip}{5pt}

\chapter{Introducción}

El objetivo de este proyecto es la predicción temprana de la enfermedad de Alzheimer a partir de imágenes biomédicas de pacientes en distintos estados de la endermedad. Este proceso de predicción se realizará utilizando técnicas que se engloban dentro del Deep Learning.

\section{Motivación}

Actualmente, la enfermadad de Alzheimer es la principal causa de demencia en el ser humano. Lo que se produce es un deterioro progresivo de las celulas nerviosas cerebrales, lo que desemboca en demencia senil. El aumento de Apersonas afectadas por esta enfermedad viene dado en parte por el aumento de la esperanza de vida, ya que esta enfermedad se da en personas de edad avanzada. Las consecuencias de esta enfermedad son varias, pero todas limitan la calidad de vida de la persona que la padece, además del de las personas de su entorno, ya que esta enfermedad va minando la autonomía del enfermo sobre su propio cuerpo. Esto provocará que el enfermo necesite un apoyo externo para poder realizar su día a día.\\
Aún no se sabe la causa de la enfermedad de Alzheimer, aunque si se conocen factores de riesgo que pueden ayudar a su desarrollo \cite{RiskFactors}. Los principales que se conocen son los siguientes: la edad (como se ha explicado antes con el aumento de la esperanza de vida), el sexo (afectando más a mujeres que a hombres), el nivel de eduación (las personas con estudios y mayor actividad intelectual son menos propensas a esta enfermedad), los trastornos metabólicos asociados con la resistencia a la insulina, la obesidad, hipertensión o diabetes y la genética (la presencia del genotipo determinado del gen de la EPOE) entre otros factores.\\
No existe una cura de la enfermedad, solo se conocen tratamientos que relentizan su avance. Estos tratamientos deben realizarse en las primeras fases de la enfermedad, ya que después pueden resultar perjudiciales para el paciente.\\
El diagnóstico del Alzheimer se podría dividir en tres partes \cite{Diagnostico}:
\begin{itemize}
	\item \textbf{Evaluación de estado de ánimo y estado mental:} El estado mental se evalúa para darle al médico una idea general de qué tan bien está funcionando la mente. Este examen da un sentido general sobre si la persona: está consciente de que tiene síntomas. sabe la fecha, la hora, y dónde está ella, puede recordar una lista corta de palabras, seguir instrucciones y hacer cálculos simples. El doctor puede que le pregunte a la persona su dirección, qué año es y quién es el presidente del país. Puede que al individuo se le pida que deletree una palabra al revés, que dibuje un reloj y que copie un diseño. El doctor también puede evaluar el estado de ánimo y el sentimiento de bienestar de la persona para detectar si hay depresión u otra enfermedad que puede causar pérdida de memoria y confusión.
	\item \textbf{Examen físico y exámenes para el diagnóstico}: El doctor hará ciertos procedimientos para evaluar la salud en general de la persona como su dieta, tomar la presión arterial, o escuchar su corazón. Se harán pruebas de sangre y de orina y posiblemente se ordenen otros exámenes de laboratorio. La información que pueden proporcionar estos exámenes puede ayudar a identificar problemas como anemia, diabetes, problemas de los riñones o del hígado, deficiencias de vitaminas, anormalidades de la tiroides y problemas del corazón o de los vasos sanguíneos. Todas estas condiciones pueden causar confusión, problemas de memoria u otros síntomas similares a la Demencia.
	\item \textbf{Examen neurológico:} Un doctor o a veces un Neurólogo que se especializa en problemas del cerebro y del sistema nervioso, evaluará muy cuidadosamente a la persona para determinar si hay señales de otro tipo de problema del cerebro que no es Alzheimer. El doctor también va a evaluar los reflejos de la persona, el equilibrio, movimiento de los ojos, lenguaje y sensibilidad. El doctor está buscando señales de derrames pequeños o grandes, enfermedad de Parkinson, tumores cerebrales, acumulación de líquido en el cerebro y otros males que pueden perjudicar la memoria o la capacidad de pensar. El examen neurológico puede incluir un MRI (Imagen por Resonancia Magnética) o CT (tomografía computarizada). Los MRI y CT pueden revelar tumores, evidencia de derrames pequeños o grandes, daño a causa de lesiones por traumas severos de cabeza o acumulación de líquido. Medicare cubre el PET (tomografía por emisión de positrones) como una ayuda para el diagnóstico en ciertos casos .
\end{itemize}
El análisis de la imagen de resonancia magnética a través del ojo humano resulta sencillo si las alteraciones estructurales son apreciables. Normalmente se utiliza una escala de grises para poder resaltar las diferencias \cite{FormatosImagenes}. El número de bits con el que trabajan estos sistemas suele ser de 16 bits \cite{FormatosImagenes} lo cuál da una escala de grises de hasta 65536 niveles, muy superior a los niveles que el ojo puede diferenciar.\\
Para tener un estudio más exhaustivo del estado del paciente, se aplican una serie de algoritmos de ordenador para poder extraer las zonas de mayor interés que son las que observa el evaluador para poder determinar si existen indicios de la enfermedad en el paciente. El principal problema es que si nos encontramos en el inicio de la enfermedad, pueden no existir una alteración suficiente para poder diagnosticar la enfermedad por el ojo humano. He aquí la motivación de la realización del proyecto, poder ayudar a un diagnóstico temprano de la enfermadad para poder tratarla con procedimientos que reduzcan su avance, y mejoren la vida del paciente.

\section{Objetivos}

El objetivo que se persigue con este proyecto es la predicción temprana de la enfermedad del Alzheimer en personas que la comienzan a padecer. El estudio que se realiza puede dividirse en los siguientes objetivos específicos:
\begin{itemize}
	\item \textbf{Conversión de las imágenes para su tratamiento:} Las imágenes de la base de datos se encuentran en formato NiFTI. Para su tratamiento con los algoritmos, debemos convertirlas a un formato con el que se puedan tratar. Con esto se creará la base de datos de la cual se aprenderá para realizar las predicciones.
	\item \textbf{Elección de la arquitectura de la Convolutional Neural Network:} se deberá elegir una arquitectura para la red neuronal.
	\item \textbf{Entrenamiento de la red neuronal:} Entrenamiento de la red neuronal con la base de datos que se ha obtenido previamente.
	\item \textbf{Análisis de los resultados de clasificación y predicción:} En este punto, se analizarán los resultados obtenidos, donde se podrá observar si las técnicas aplicadas tienen un buen rendimiento.
\end{itemize}
\section{Resumen y estructura del proyecto}
\bibliography{bibliografia/bibliografia.bib}
\bibliographystyle{ieeetr} % hay varias formas de citar

%\input{capitulos/02_EspecificacionRequisitos}
%
%\input{capitulos/03_Planificacion}
%
%\input{capitulos/04_Analisis}
%
%\input{capitulos/05_Diseno}
%
%\input{capitulos/06_Implementacion}
%
%\input{capitulos/07_Pruebas}
%
%\mychapter{8}{Conclusiones y trabajos futuros}

En este capítulo se va a hablar de las conclusiones que se pueden observar a partir de los resultados obtenidos en el capítulo anterior y trabajos futuros que se desearían realizar a partir de lo ya concluido con este trabajo.\\

\section{Conclusiones}

Con este trabajo se intentaban dar respuesta a varias cuestiones. Una sería si \textbf{se pueden obtener buenos resultados usando técnicas de \textit{Deep Learning} para la clasificación de imágenes de pacientes con Alzheimer y pacientes sin esta enfermedad usando imágenes en 2D}, al contrario que en estudios previos donde se utilizaban las imágenes en 3D \cite{residualVGG}.\\

Como se puedo observar en los experimentos previos \ref{experimentosprevios-resultados} y con el experimento 2 \ref{experimento2-resultados}, los resultados mejoran cuanto mayor es la base de datos que tenemos. Esto deja dilucidar que si aumentamos la base de datos cada vez con más imágenes se podrían obtener muy buenos resultados a la hora de la clasificación de pacientes que nos llegasen, pudiendo llegar a ser una herramienta de utilidad a la hora de que un médico pudiese hacer un diagnóstico a un paciente. Además, al tratarse de imágenes 2D en lugar de 3D, el tiempo para el entrenamiento y la capacidad de computación necesaria para el mismo o la predicción se reducen notablemente, dejando un camino abierto para que se pudiese realizar en un ordenador menos potente, lo cuál abriría más caminos para que se pudiese implantar un sistema así en un hospital.\\

La segunda cuestión a la que se intentaba dar respuesta era que, una vez visto que la clasificación de los pacientes con imágenes 2D es viable y se pueden observar mejoras, \textbf{cuál es la capa del cerebro con la que se debería hacer esta clasificación}. Para ello se preparó el experimento 3, y se pueden ver sus resultados en \ref{experimento3-resultados}. Con estos en mano podemos observar que las mejores capas para la clasificación son la \textbf{40,55,50}. Aun así, son un poco optimistas, ya que vimos en el experimento 2 \ref{experimento2-resultados} que nuestro techo de \textit{Accuracy} utilizando un método de validación muy robusto como es el \textit{LOO CV} era de un \textbf{60'5655\%}. Con estos resultados también hemos podido afirmar que hay unas capas que funcionan mejor que otras.\\

Además, como también hemos comprobado que hay una filtración de información entre distintas imágenes del mismo paciente con los resultados en \ref{experimentosprevios-resultados} y \ref{experimento1-resultados}, siempre tenemos que realizar una separación de los pacientes y esto se debería realizar también a la hora de que se implantase este procedimiento en un centro médico. \\

Esto por esto que los resultados \textbf{más significativos y en los que se basa este trabajo} son los obtenidos en los experimentos 2 \ref{experimento2-resultados} y 3 \ref{experimento3-resultados}.\\

Gracias a estos resultados podríamos afirmar que hay un futuro a la hora de realizar predicciones sin utilizar imágenes 3D, reduciendo los costes de computación y sin realizar un procesamiento de las imágenes muy costoso como se ha explicado en \ref{procesamiento-imagenes}, pero que se debe disponer de una mayor base de datos para que incrementemos la generalización de la red. Además, se ha confirmado que la estructura de la red utilizada en \cite{residualVGG} es útil también para la clasificación de imágenes 2D, obteniendo unos resultados aceptables en \ref{experimento2-resultados} y \ref{experimento3-resultados}.\\

Lo que también se puede afirmar es que un mayor número de imágenes de pacientes da lugar a unos mejores resultados, como se pudo observar en los experimentos previos que se realizaron \ref{experimentosprevios-resultados}. Esto indicaría que si se pudiese normalizar la base de datos actual y aplicarle el post procesamiento \ref{postproces} al que fue sometida la base de datos utilizada para este trabajo, se podrían obtener unos muy buenos resultados a la hora de la clasificación de los pacientes.\\

\section{Trabajos futuros}

En cuanto a los trabajos futuros a partir del actual, el principal sería continuar por la línea de comprobar si cuantas más imágenes mejores resultados. Actualmente solo se ha utilizado un subgrupo dentro de la base de datos donde se había aplicado el post procesamiento explicado en \ref{postproces}. Como continuación, nos gustaría aplicar el post procesamiento a un número mayor de imágenes y obtener nuevos resultados.\\

Además, después de haber respondido a la pregunta de que capas funcionan mejor que otras, lo siguiente sería responder porque unas capas funcionan mejores que otras. Este es otro estudio muy interesante el cuál sería recomendable realizar para darle una respuesta a estas preguntas.\\

Como último, una continuación muy interesante una vez que se haya respondido a estas cuestiones y se hayan conseguido buenos resultados, sería la creación de una plataforma web donde un doctor pudiese subir las imágenes de un paciente y se produjese una predicción que el doctor pudiese tomar como primera consulta a la hora de diagnosticar al paciente. Esto sería realizable una vez que se hubiesen obtenido buenos resultados y estos tuviesen una robustez muy considerable.\\
%
%%\chapter{Conclusiones y Trabajos Futuros}
%
%
%\nocite{*}
\bibliography{bibliografia/bibliografia.bib}
\bibliographystyle{ieeetr}
\addcontentsline{toc}{chapter}{Bibliografía}
%\bibliographystyle{miunsrturl}

%
%\appendix
%\input{apendices/manual_usuario/manual_usuario}
%%\input{apendices/paper/paper}
%\input{glosario/entradas_glosario}
% \addcontentsline{toc}{chapter}{Glosario}
% \printglossary
\chapter*{}
\thispagestyle{empty}

\end{document}
